\section{Cancellation carried by the true section source}
\label{sec:v39-section-source}
The existence of a peripheral phase is not by itself a nonzero
contribution to a prescribed endpoint pairing. For the actual section
we can calculate a source cancellation on the full occupation torus.
The calculation takes place in the physical $L^2$ space and requires
neither a full-torus anisotropic power bound nor regularity of an
eigenfunction.

Use the inner product $\langle f,g\rangle=\int\overline f g\,d\nu$.
Let $f_R^{\rm c}=(\kappa_{R,1},\kappa_{R,2},\tau_R-\bar\tau_R)$,
$z=e^{iv}\ne1$, and define the physical unitary operator
\begin{equation}\label{eq:v39-physical-unitary}
 U_{R,\omega,v}h
  =\bigl(e^{i(\omega\cdot f_R^{\rm c}+v\eta_R)}h\bigr)\circ T_R^{-1}.
\end{equation}
It is the physical realization of $e^{ivc}\mathcal Q_{R,\omega,v}$.
Unlike the anisotropic realization, its full $L^2$ operator norm is
one at every power. Its inverse is
$U^{-1}h=e^{-i(\omega\cdot f_R^{\rm c}+v\eta_R)}h\circ T_R$.

\subsection{A source resolvent, not a full-space resolvent bound}
\begin{theorem}[Section-source spectral cancellation]
\label{thm:v39-section-source}
For $|\omega|\le1$, $v\notin2\pi\Z$, $0\le r<1$ and
$0<\epsilon\le2$, let $E_{R,\omega,v}(\epsilon)$ be the spectral
projection of $U_{R,\omega,v}$ on $|\lambda-1|\le\epsilon$.
Uniformly in $R$,
\begin{align}
 \|E_{R,\omega,v}(\epsilon)\eta_R\|_2
 &\le\frac{\epsilon+C|\omega|}{|e^{iv}-1|},
                                  \label{eq:v39-source-spectral-window}\\
 \|(I-rU_{R,\omega,v})^{-1}\eta_R\|_2
 &\le\frac{2/(1+r)+C|\omega|/(1-r)}{|e^{iv}-1|},
                                  \label{eq:v39-source-resolvent}\\
 \left\|\frac1M\sum_{j=0}^{M-1}U_{R,\omega,v}^j\eta_R\right\|_2
 &\le\frac{2/M+C|\omega|}{|e^{iv}-1|}\quad(M\ge1).
                                  \label{eq:v39-source-cesaro}
\end{align}
At $\omega=0$, the spectral measures of $\eta_R$ and $1$ satisfy the
exact identity
\begin{equation}\label{eq:v39-source-spectral-measure}
 d\mu_{\eta_R,v}(\lambda)
   =\frac{|\lambda-1|^2}{|e^{iv}-1|^2}\,d\mu_{1,v}(\lambda).
\end{equation}
In particular the projection of $\eta_R$ onto the eigenvalue $1$
is zero at every nonzero occupation phase, including an occupation
phase for which that eigenvalue exists.
\end{theorem}
\begin{proof}
The elementary identity $e^{-iv\eta_R}=1+(z^{-1}-1)\eta_R$ and
the formula for $U^{-1}$ give
\begin{equation}\label{eq:v39-source-identity}
 (z^{-1}-1)\eta_R=(U^{-1}-I)1+h_{R,\omega,v},\qquad
 h_{R,\omega,v}=(e^{i\omega\cdot f_R^{\rm c}}-1)U^{-1}1.
\end{equation}
The bounded one-collision record implies
$\|h_{R,\omega,v}\|_2\le C|\omega|$. This estimate is independent
of $v$; no derivative of an occupation itinerary occurs.

Apply the spectral projection in the statement to
\eqref{eq:v39-source-identity}. On its range, $U^{-1}-I$ has norm
at most $\epsilon$. The projection is a contraction and $\|1\|_2=1$.
Division by $|z^{-1}-1|=|z-1|$ proves
\eqref{eq:v39-source-spectral-window}.
For the resolvent, the spectral theorem gives
\[
 \sup_{|\lambda|=1}\frac{|\lambda^{-1}-1|}{|1-r\lambda|}
       =\frac2{1+r},\qquad
 \|(I-rU)^{-1}\|\le\frac1{1-r}.
\]
Indeed $|1-r\lambda|^2=(1-r)^2+r|\lambda-1|^2$, and the ratio
is increasing in $|\lambda-1|^2\in[0,4]$. Apply these bounds to the
two terms in \eqref{eq:v39-source-identity}. For the Cesaro bound,
\[
 \sum_{j=0}^{M-1}U^j(U^{-1}-I)1=(U^{-1}-U^{M-1})1
\]
has norm at most two, and the average applied to $h$ has norm at most
$\|h\|_2$. This proves \eqref{eq:v39-source-cesaro}.
At $\omega=0$, $h=0$ in the exact source identity. Functional calculus
then proves \eqref{eq:v39-source-spectral-measure}, including its
assertion at the atom $\lambda=1$.
\end{proof}

\begin{corollary}[An invisible pure occupation eigenphase]
\label{cor:v39-invisible-occupation-phase}
If $q$ is a measurable circle-valued solution of
$q\circ T_R=e^{iv\eta_R}q$ with $v\notin2\pi\Z$, then
\[
 \int\eta_Rq\,d\nu=0,\qquad \int q\,d\nu_R^*=0.
\]
Thus its rank-one physical eigenspace has zero coefficient in the
$\eta_R$--$\eta_R$ pairing. No analogous vanishing is asserted for
an arbitrary physical frequency or for an eigenvalue different from one.
\end{corollary}
\begin{proof}
Invariance gives
$0=\int(q\circ T_R-q)\,d\nu=(e^{iv}-1)\int\eta_Rq\,d\nu$.
Division by $c$ proves the second equality. Equivalently,
Proposition~\ref{prop:v39-induced-arithmetic} restricts $q$ to an
induced eigenfunction with eigenvalue $e^{iv}$, whose mean is zero by
measure preservation. The physical eigenspace at eigenvalue one is at
most one dimensional: the quotient of two circle solutions is invariant.
The assertion about the pairing follows, also from
Theorem~\ref{thm:v39-section-source}.
\end{proof}

\subsection{An exact finite-time cancellation and its Abel limit}
For $m\ge0$ put
\[
 F_{m,R}(v)=\int e^{ivA_{m,R}}\,d\nu,\qquad F_{0,R}=1,
\]
and for $m\ge1$ put
\begin{equation}\label{eq:v39-section-characteristic}
 C_{m,R}(v)=\int\eta_R(x)\eta_R(T_R^mx)e^{ivA_{m,R}(x)}\,d\nu(x).
\end{equation}
By the exact return disintegration, the same function is
\begin{equation}\label{eq:v39-return-generating-coefficient}
 C_{m,R}(v)=c\sum_{n\ge1}e^{ivn}\nu_R^*\{N_{n,R}=m\}.
\end{equation}
Only finitely many terms occur at any given $m$.

\begin{theorem}[Full-torus section cancellation in the collision clock]
\label{thm:v39-abel-section}
For every $m\ge1$ and $v\notin2\pi\Z$,
\begin{equation}\label{eq:v39-second-difference}
 C_{m,R}(v)=\frac{z}{(z-1)^2}
       \bigl(F_{m+1,R}(v)-2F_{m,R}(v)+F_{m-1,R}(v)\bigr).
\end{equation}
Consequently, uniformly in $R$ and $M\ge1$,
\begin{equation}\label{eq:v39-partial-sum-cancellation}
 \left|\sum_{m=1}^M C_{m,R}(v)\right|
                         \le\frac4{|e^{iv}-1|^2}.
\end{equation}
For $1/2\le r<1$ the absolutely convergent Abel sum obeys
\begin{equation}\label{eq:v39-abel-limit}
 \left|\sum_{m\ge1}r^m C_{m,R}(v)-\frac{cz}{1-z}\right|
                 \le\frac{4(1-r)}{|e^{iv}-1|^2}.
\end{equation}
In particular these estimates are uniform on the entire occupation
torus outside any fixed neighborhood of zero.
\end{theorem}
\begin{proof}
Write $\eta_j=\eta_R\circ T_R^j$. Stationarity gives
\[
 F_{m+1,R}-F_{m,R}
  =(z-1)\int\eta_0 z^{\eta_1+\cdots+\eta_m}\,d\nu.
\]
Subtract its version with $m$ replaced by $m-1$. The result is
$(z-1)^2\int\eta_0\eta_m z^{\eta_1+\cdots+\eta_{m-1}}\,d\nu$.
On the support of $\eta_0\eta_m$, multiplication by $z$ inserts the
initial visit and not the terminal one. This is precisely
\eqref{eq:v39-second-difference}. The convention also covers $m=1$.
Summing the second difference leaves
$F_{M+1,R}-F_{M,R}-F_{1,R}+F_{0,R}$, whose modulus is at most four.
This proves \eqref{eq:v39-partial-sum-cancellation}.

Let $F_R(r,v)=\sum_{m\ge0}r^m F_{m,R}(v)$. Absolute convergence
follows from $|F_{m,R}|\le1$. Since $F_{1,R}=1+c(z-1)$, summing
\eqref{eq:v39-second-difference} gives the exact identity
\begin{equation}\label{eq:v39-exact-abel}
 \sum_{m\ge1}r^m C_{m,R}(v)
 =\frac{z}{(z-1)^2}\left\{
   \frac{(1-r)^2}{r}F_R(r,v)-\frac{1-r}{r}-c(z-1)\right\}.
\end{equation}
The first two terms in braces have total modulus at most
$2(1-r)/r\le4(1-r)$. The last term gives $cz/(1-z)$, proving
\eqref{eq:v39-abel-limit}. No mixing estimate was used: the cancellation
comes from the same initial and terminal section indicator and from
integer-valued occupation.
\end{proof}

\paragraph{Relation to the fixed-count inverse.}
These are new estimates for the unmodified section source. They show
that a pure occupation eigenphase at eigenvalue one is invisible to the
required endpoints, and that a whole occupation-torus complement has
bounded collision-clock partial sums and an explicit Abel limit.
They do not give decay of $C_{m,R}(v)$ for each fixed $m$.
For example, on a stationary two-cycle with section indicator $(1,0)$
and $z=-1$, one has $C_{2j}(v)=(-1)^j/2$ and $C_{2j+1}(v)=0$.
Its partial sums are bounded but its coefficients do not tend to zero.
This example is not the billiard and is used only to distinguish the
two conclusions of the source identity.

In particular the $m^2$-normalized remainder of
Proposition~\ref{prop:v38-exact-index-remainder} is still a fixed-$m$
coefficient and includes nonzero displacement--roof frequencies. Neither
its cancellation nor a pointwise roof-density estimate follows by
removing the collision-clock average from \eqref{eq:v39-abel-limit}.
The original raw-return correction, full complement and weighted
microscopic targets remain unchanged.
